% ============================================================
% Paper 88 (DE): Ungerade vollkommene Zahlen:
%                Struktur, Schranken und das offene Problem
% Autor: Michael Fuhrmann
% Specialist Mathematics Project, Build 168
% Datum: 2026-03-12
% ============================================================
\documentclass[12pt,a4paper]{amsart}

\usepackage{amsmath,amssymb,amsthm,mathtools,enumitem,booktabs,array,hyperref}
\usepackage[utf8]{inputenc}
\usepackage[T1]{fontenc}
\usepackage{geometry}
\geometry{margin=2.5cm}

% ---- Theorem-Umgebungen (Deutsch) ----
\newtheorem{theorem}{Theorem}[section]
\newtheorem{lemma}[theorem]{Lemma}
\newtheorem{corollary}[theorem]{Korollar}
\newtheorem{proposition}[theorem]{Proposition}
\newtheorem{satz}[theorem]{Satz}
\newtheorem{korollar}[theorem]{Korollar}
\theoremstyle{definition}
\newtheorem{definition}[theorem]{Definition}
\newtheorem{definition_de}[theorem]{Definition}
\newtheorem{example}[theorem]{Beispiel}
\theoremstyle{remark}
\newtheorem{remark}[theorem]{Bemerkung}
\newtheorem{bemerkung}[theorem]{Bemerkung}
\newtheorem{conjecture}[theorem]{Vermutung}
\newtheorem{vermutung}[theorem]{Vermutung}

% ---- Operatoren ----
\DeclareMathOperator{\ord}{ord}
\newcommand{\N}{\mathbb{N}}
\newcommand{\Z}{\mathbb{Z}}
\newcommand{\Q}{\mathbb{Q}}
\newcommand{\R}{\mathbb{R}}

% ============================================================
\begin{document}

\title{Ungerade vollkommene Zahlen:\\
Struktur, Schranken und das offene Problem}

\author{Michael Fuhrmann}
\address{Specialist Mathematics Project, Build 168}
\email{specialist-maths@project.local}
\date{2026-03-12}

\begin{abstract}
Eine positive ganze Zahl $n$ heißt \emph{vollkommen}, wenn $\sigma(n) = 2n$ gilt,
wobei $\sigma$ die Summe aller Teiler bezeichnet.
Alle bekannten vollkommenen Zahlen sind gerade und entstehen aus Mersenne-Primzahlen
gemäß dem Satz von Euklid und Euler.
Die Frage, ob eine \emph{ungerade} vollkommene Zahl existiert, ist seit über
zweitausend Jahren ungelöst und gilt als das älteste offene Problem der Mathematik.
Diese Arbeit gibt einen Überblick über die wichtigsten strukturellen Schranken,
die eine hypothetische ungerade vollkommene Zahl erfüllen müsste:
Eulers Faktorisierungssatz ($n = p^a m^2$ mit $p \equiv a \equiv 1 \pmod{4}$),
moderne untere Schranken ($n > 10^{1500}$, Ochem--Rao 2012),
Schranken für die Anzahl verschiedener Primteiler ($\omega(n) \geq 9$),
Schranken für den größten und kleinsten Primteiler sowie
Teilbarkeits- und Kongruenzbeschränkungen.
Weiterhin werden verwandte Konzepte --- quasivollkommene, mehrfach vollkommene
und hypervollkommene Zahlen --- diskutiert.
\end{abstract}

\maketitle

\tableofcontents

% ============================================================
\section{Einleitung}
% ============================================================

Die \emph{Teilersummenfunktion} $\sigma : \N \to \N$ ordnet jeder positiven ganzen
Zahl $n$ die Summe aller ihrer positiven Teiler zu:
\[
  \sigma(n) = \sum_{d \mid n} d.
\]
Eine Zahl heißt \emph{vollkommen}, wenn $\sigma(n) = 2n$, d.\,h.\ wenn die echten
Teiler von $n$ genau $n$ ergeben.
Die ersten vollkommenen Zahlen sind $6 = 1+2+3$, $28 = 1+2+4+7+14$,
$496$ und $8128$.

Die antiken Griechen, insbesondere Euklid und Nikomachus, waren von vollkommenen
Zahlen fasziniert.
Euklid bewies, dass $2^{p-1}(2^p-1)$ vollkommen ist, wenn $2^p - 1$ prim ist
(eine \emph{Mersenne-Primzahl}).
Euler zeigte später, dass \emph{jede} gerade vollkommene Zahl diese Form hat.
Damit sind gerade vollkommene Zahlen vollständig klassifiziert --- vorbehaltlich
der offenen Frage nach der Anzahl der Mersenne-Primzahlen.

Die Situation für \emph{ungerade} vollkommene Zahlen ist grundlegend verschieden.
Trotz jahrhundertelanger Bemühungen und moderner Computersuchen, die
$n > 10^{1500}$ erreichen, wurde noch keine ungerade vollkommene Zahl gefunden,
und niemand hat bewiesen, dass keine existiert.

\begin{vermutung}[Keine ungerade vollkommene Zahl]
\label{verm:keine}
Es gibt keine ungerade vollkommene Zahl.
\end{vermutung}

Vermutung~\ref{verm:keine} gilt als weitgehend richtig.
Diese Arbeit stellt die stärksten bekannten strukturellen Bedingungen vor,
die eine hypothetische ungerade vollkommene Zahl erfüllen müsste.

% ============================================================
\section{Arithmetischer Hintergrund}
% ============================================================

\subsection{Die Teilersummenfunktion}

\begin{definition_de}
Für $n \in \N$ sei
\[
  \sigma(n) = \sum_{d \mid n} d, \qquad
  s(n) = \sigma(n) - n \quad \text{(Summe der \emph{echten} Teiler)}.
\]
\end{definition_de}

\begin{satz}[Multiplikativität]
$\sigma$ ist multiplikativ auf teilerfremden ganzen Zahlen und erfüllt auf Primzahlpotenzen:
\[
  \sigma(p^k) = 1 + p + \cdots + p^k = \frac{p^{k+1}-1}{p-1},
\]
und für $\gcd(a,b)=1$: $\sigma(ab) = \sigma(a)\sigma(b)$.
\end{satz}

\begin{proof}
Jeder Teiler von $p_1^{k_1}\cdots p_r^{k_r}$ lässt sich eindeutig als
$\prod p_i^{e_i}$ mit $0 \leq e_i \leq k_i$ schreiben; die Summe
faktorisiert entsprechend.
\end{proof}

\begin{definition_de}
Der \emph{Abundanzindex} von $n$ ist $\sigma(n)/n$.
Eine vollkommene Zahl erfüllt $\sigma(n)/n = 2$.
\end{definition_de}

\subsection{Klassifikation nach Abundanz}

\begin{definition_de}
Eine positive ganze Zahl $n$ heißt:
\begin{itemize}[nosep]
  \item \emph{defizient}, wenn $\sigma(n) < 2n$;
  \item \emph{vollkommen}, wenn $\sigma(n) = 2n$;
  \item \emph{abundant}, wenn $\sigma(n) > 2n$.
\end{itemize}
\end{definition_de}

Jede Primzahl $p$ ist defizient: $\sigma(p)/p = 1 + 1/p < 2$.
Potenzen von $2$ erfüllen $\sigma(2^k)/2^k = 2 - 2^{1-k} < 2$, sind also
ebenfalls defizient.

% ============================================================
\section{Eulers Struktursatz}
% ============================================================

Das grundlegende Resultat in der Theorie ungerader vollkommener Zahlen
stammt von Euler:

\begin{satz}[Eulers Struktursatz, 1747]
\label{satz:euler}
Ist $n$ eine ungerade vollkommene Zahl, so gilt
\[
  n = p^a \cdot m^2,
\]
wobei $p$ eine Primzahl, $\gcd(p, m) = 1$, $p \equiv 1 \pmod{4}$
und $a \equiv 1 \pmod{4}$.
\end{satz}

\begin{proof}[Beweisidee]
Schreibe $n = p_1^{a_1} \cdots p_r^{a_r}$ mit allen $p_i$ ungerade.
Da $\sigma(n) = 2n$, gilt
\[
  \prod_{i=1}^r \sigma(p_i^{a_i}) = 2 \prod_{i=1}^r p_i^{a_i}.
\]
Die linke Seite muss genau einen Faktor $2$ enthalten (da $n$ ungerade ist
und $2n$ genau einen zusätzlichen Faktor $2$ gegenüber $n$ hat).
Nun ist $\sigma(p^a) = 1 + p + \cdots + p^a$ für ungerades $p$ genau dann
ungerade, wenn $a$ gerade ist (da dann $a+1$ Summanden alle ungerade sind
und ihre Anzahl gerade, also $\sigma(p^a) \equiv 0 \pmod{2}$ nur wenn
$a+1$ ungerade, d.\,h.\ $a$ gerade).

Daher muss genau ein Exponent $a_j$ ungerade sein; nenne diese Primzahl
$p = p_j$ und ihren Exponenten $a = a_j$.
Alle anderen $a_i$ sind gerade, sodass $n = p^a \cdot m^2$.

Für $p \equiv 3 \pmod{4}$ und $a \equiv 1 \pmod{4}$ folgt durch
$\pmod{4}$-Rechnung, dass $\sigma(p^a)$ gerade wäre --- Widerspruch.
Also $p \equiv 1 \pmod{4}$ und $a \equiv 1 \pmod{4}$.
\end{proof}

\begin{bemerkung}
Die Primzahl $p$ im Eulerschen Struktursatz heißt die \emph{Euler-Primzahl}
oder \emph{spezielle Primzahl} von $n$.
Sie erfüllt $p \equiv 1 \pmod{4}$ und erscheint in einer ungeraden Potenz
$a \equiv 1 \pmod{4}$.
Die kleinste Möglichkeit ist $a = 1$, $p = 5$.
\end{bemerkung}

\begin{korollar}
\label{kor:form}
Jede ungerade vollkommene Zahl erfüllt $n \equiv 1 \pmod{4}$.
\end{korollar}

\begin{proof}
$n = p^a m^2$ mit $p \equiv 1 \pmod{4}$ und $m$ ungerade,
also $m^2 \equiv 1 \pmod{4}$ und $p^a \equiv 1 \pmod{4}$,
daher $n \equiv 1 \pmod{4}$.
\end{proof}

% ============================================================
\section{Schranken für die Größe}
% ============================================================

\begin{satz}[Ochem--Rao, 2012]
\label{satz:ochem}
Ist $n$ eine ungerade vollkommene Zahl, so gilt $n > 10^{1500}$.
\end{satz}

Dieses bemerkenswerte Ergebnis wurde von Ochem und Rao durch eine Kombination
algebraischer Einschränkungen und computergestützter Fallunterscheidungen erzielt.
Frühere Schranken entwickelten sich wie folgt:
\begin{itemize}[nosep]
  \item Kanold (1941): $n > 10^{20}$
  \item Brent--Cohen--te Riele (1991): $n > 10^{300}$
  \item Ochem--Rao (2012): $n > 10^{1500}$
\end{itemize}

\begin{bemerkung}
Die Schranke $n > 10^{1500}$ bedeutet, dass eine ungerade vollkommene Zahl
mindestens $1500$ Dezimalstellen hätte --- weit jenseits jeder rechnerischen
Überprüfbarkeit durch direkte Suche.
\end{bemerkung}

% ============================================================
\section{Schranken für die Anzahl der Primteiler}
% ============================================================

\begin{satz}[Chein, Hagis; mehrere Autoren]
\label{satz:omega}
Ist $n$ eine ungerade vollkommene Zahl, so gilt $\omega(n) \geq 9$,
wobei $\omega(n)$ die Anzahl der verschiedenen Primteiler von $n$ bezeichnet.
\end{satz}

Der Beweis zeigt, dass für $\omega(n) \leq 8$ die Abundanzschranke
$\sigma(n)/n = 2$ nicht erfüllt werden kann.

\begin{satz}[Chein 1979, Hagis 1980]
Gilt zusätzlich $3 \nmid n$, so folgt $\omega(n) \geq 11$.
\end{satz}

\begin{satz}[Nielsen 2015]
$\omega(n) \geq 10$, wenn die größte Komponente $p^a > n^{1/3}$.
\end{satz}

\subsection{Primteiler mit Vielfachheit}

\begin{satz}[Ochem--Rao 2014]
Ist $n$ eine ungerade vollkommene Zahl, so gilt $\Omega(n) \geq 101$,
wobei $\Omega(n)$ die Anzahl der Primteiler von $n$ mit Vielfachheit zählt.
\end{satz}

% ============================================================
\section{Schranken für einzelne Primteiler}
% ============================================================

\subsection{Größter Primteiler}

\begin{satz}[Goto--Ohno 2008]
Der größte Primteiler einer ungeraden vollkommenen Zahl übersteigt $10^8$.
\end{satz}

\begin{satz}[Luca--Pomerance 2010]
Der zweitgrößte Primteiler einer ungeraden vollkommenen Zahl übersteigt $10^4$.
\end{satz}

\subsection{Kleinster Primteiler}

\begin{satz}[Grün 1952, erweitert]
Ist $3 \mid n$, so ist der kleinste Primteiler $3$.
Ist $3 \nmid n$, so ist der kleinste Primteiler mindestens $5$.
\end{satz}

\begin{satz}[Iannucci 1999]
Ist $n$ eine ungerade vollkommene Zahl mit $3 \nmid n$ und $5 \nmid n$,
so ist der kleinste Primteiler von $n$ mindestens $100003$.
\end{satz}

% ============================================================
\section{Teilbarkeitseinschränkungen}
% ============================================================

\begin{satz}[Verschiedene Autoren]
Ist $n$ eine ungerade vollkommene Zahl, so gilt:
\begin{enumerate}[label=(\roman*)]
  \item $n \equiv 1 \pmod{4}$ (Korollar~\ref{kor:form});
  \item $n \equiv 1 \pmod{12}$ oder $n \equiv 9 \pmod{36}$ (Touchard 1953);
  \item $105 = 3 \cdot 5 \cdot 7 \nmid n$ (Steuerwald 1937);
  \item Ist $3 \mid n$, dann gilt $9 \mid n$;
  \item $n$ ist kein vollkommenes Quadrat
        (direkt aus Eulers Satz: $p^a$ mit ungeradem $a$ verhindert dies).
\end{enumerate}
\end{satz}

\begin{satz}[Touchard 1953]
\label{satz:touchard}
Ist $n$ eine ungerade vollkommene Zahl, so gilt entweder
$n \equiv 1 \pmod{12}$ oder $n \equiv 9 \pmod{36}$.
\end{satz}

\begin{proof}[Beweisidee]
Dies folgt aus Eulers Struktursatz kombiniert mit einer sorgfältigen Analyse
von $\sigma(n)/n = 2$ modulo $3$ und $4$ gleichzeitig.
\end{proof}

% ============================================================
\section{Abundanzindex und Obstruktionstheorie}
% ============================================================

\begin{definition_de}
Eine \emph{Abundanzobstruktion} für eine rationale Zahl $r$ ist ein Beweis,
dass kein $n \in \N$ die Gleichung $\sigma(n)/n = r$ erfüllt.
\end{definition_de}

\begin{satz}[Broughan--Delbourgo--Zhou]
Es gibt eine dichte Menge rationaler Zahlen $r > 1$, für die keine ganze
Zahl $n$ mit $\sigma(n)/n = r$ existiert. Die Zahl $2$ gehört nach
aktuellen Methoden jedoch nicht zu dieser Menge.
\end{satz}

\begin{bemerkung}
Die Nichtexistenz ungerader vollkommener Zahlen kann allein durch lokale
Obstruktionen nicht bewiesen werden: Für jede endliche Primzahlmenge $S$
gibt es ungerade (nicht vollkommene) Zahlen $n$, deren $S$-lokaler
Abundanzindex mit dem von $2$ übereinstimmt.
\end{bemerkung}

% ============================================================
\section{Verwandte Konzepte}
% ============================================================

\subsection{Quasivollkommene Zahlen}

\begin{definition_de}
Eine positive ganze Zahl $n$ heißt \emph{quasivollkommen}, wenn $\sigma(n) = 2n + 1$.
\end{definition_de}

\begin{vermutung}
Es gibt keine quasivollkommenen Zahlen.
\end{vermutung}

Es ist bekannt, dass jede quasivollkommene Zahl ein ungerades vollkommenes
Quadrat $> 10^{35}$ mit mindestens sieben verschiedenen Primteilern sein müsste.

\subsection{Mehrfach vollkommene Zahlen}

\begin{definition_de}
Eine positive ganze Zahl $n$ heißt \emph{$k$-fach vollkommen}, wenn
$\sigma(n) = kn$.
\end{definition_de}

Für $k = 3$ (dreifach vollkommen) sind sechs Zahlen bekannt:
$120, 672, 523776, 459818240, 1476304896, 51001180160$.
Für $k \geq 4$ sind viele Beispiele bekannt, aber alle sind gerade.

\begin{vermutung}
Für alle $k \geq 3$ sind alle $k$-fach vollkommenen Zahlen gerade.
\end{vermutung}

\subsection{Hypervollkommene Zahlen}

\begin{definition_de}
Eine positive ganze Zahl $n$ heißt \emph{$k$-hypervollkommen}, wenn
$\sigma(n) = (1+1/k)n + (1-1/k)$.
\end{definition_de}

\begin{bemerkung}
Für $k = 1$ reduziert sich hypervollkommen auf vollkommen.
Ob für $k \geq 2$ ungerade hypervollkommene Zahlen existieren, ist offen.
\end{bemerkung}

% ============================================================
\section{Verbindungen zu anderen offenen Problemen}
% ============================================================

\subsection{Verbindung zu Lehmer-Zahlen}

Ungerade vollkommene Zahlen stehen in Verbindung mit Lehmers Totient-Problem:
Falls $\phi(n) \mid (n-1)$ für zusammengesetztes $n$ gilt (eine Lehmer-Zahl),
ähneln die strukturellen Einschränkungen denen ungerader vollkommener Zahlen.

\subsection{Unitär vollkommene Zahlen}

\begin{definition_de}
$n$ heißt \emph{unitär vollkommen}, wenn $\sigma^*(n) = 2n$, wobei
$\sigma^*(n) = \sum_{d \mid n,\, \gcd(d,n/d)=1} d$.
\end{definition_de}

Nur fünf unitär vollkommene Zahlen sind bekannt; alle sind gerade.

% ============================================================
\section{Computergestützte Suche}
% ============================================================

\begin{satz}[Ochem--Rao 2012, rechnerisch]
Es gibt keine ungerade vollkommene Zahl kleiner als $10^{1500}$.
\end{satz}

Moderne Computersuchen nutzen:
\begin{enumerate}[label=(\arabic*)]
  \item \textbf{Faktorkettenverfahren}: Erzwingung von $\sigma(p^a) \mid 2n$
        für jede Primzahlpotenzkomponente.
  \item \textbf{Erschöpfende Fallanalyse}: nach Anzahl der Primteiler,
        Größe der Euler-Primzahl usw.
  \item \textbf{Algebraische Schranken}: Kombination aus Touchards Satz,
        Kongruenzbedingungen und Eulerscher Faktorisierung.
\end{enumerate}

\begin{bemerkung}
Die rechnerische Aufgabe, alle ungeraden Zahlen bis $10^{100}$ direkt auf
Vollkommenheit zu prüfen, ist völlig undurchführbar; die Schranken werden
durch strukturelle Argumente, nicht durch erschöpfende Suche erzielt.
\end{bemerkung}

% ============================================================
\section{Heuristische Argumente}
% ============================================================

\subsection{Wahrscheinlichkeitsheuristik}

Eine Standardheuristik legt nahe, dass die Anzahl vollkommener Zahlen
bis $x$ sehr langsam wächst.
Für ungerade Zahlen argumentiert man wie folgt:
Die Wahrscheinlichkeit, dass eine ``zufällige'' ungerade Zahl $n$ die
Gleichung $\sigma(n)/n = 2$ erfüllt, hängt von der Dichte der Menge
$\{r \in \Q : r = \sigma(n)/n\}$ nahe bei $2$ ab.

\begin{bemerkung}[Heuristisch]
Da $\sigma(n)/n$ in strukturierten Familien variiert und die Einschränkung
äußerst rigid ist (exakt $2$, nicht näherungsweise), legen heuristische
probabilistische Argumente nahe, dass die erwartete Anzahl ungerader
vollkommener Zahlen $0$ beträgt, was mit Vermutung~\ref{verm:keine}
übereinstimmt.
\end{bemerkung}

% ============================================================
\section{Zusammenfassung notwendiger Bedingungen}
% ============================================================

Wir sammeln alle etablierten notwendigen Bedingungen:

\begin{satz}[Notwendige Bedingungen für ungerade vollkommene Zahlen]
\label{satz:nw}
Existiert eine ungerade vollkommene Zahl $n$, so gilt:
\begin{enumerate}[label=(\roman*)]
  \item $n = p^a m^2$ mit Primzahl $p$, $p \equiv a \equiv 1 \pmod{4}$,
        $\gcd(p,m)=1$;
  \item $n > 10^{1500}$;
  \item $\omega(n) \geq 9$;
  \item $\Omega(n) \geq 101$;
  \item Der größte Primteiler von $n$ übersteigt $10^8$;
  \item Der zweitgrößte Primteiler übersteigt $10^4$;
  \item $n \equiv 1 \pmod{4}$;
  \item $n \equiv 1 \pmod{12}$ oder $n \equiv 9 \pmod{36}$;
  \item $105 \nmid n$;
  \item $n$ ist kein vollkommenes Quadrat.
\end{enumerate}
\end{satz}

\begin{bemerkung}
Satz~\ref{satz:nw} fasst die angesammelten Erkenntnisse von über zwei
Jahrhunderten Zahlentheorie zusammen.
Jede Bedingung schränkt den Suchraum erheblich ein --- zusammen liefern sie
jedoch noch keinen Widerspruch.
\end{bemerkung}

% ============================================================
\section{Schluss und offene Fragen}
% ============================================================

Das Problem ungerader vollkommener Zahlen ist eines der ältesten offenen
Probleme der Mathematik, ohne Aussicht auf baldige Lösung trotz des
kombinierten Einsatzes analytischer, algebraischer und rechnerischer Methoden.

Die zentrale Vermutung bleibt:

\begin{vermutung}[Wiederholt]
Es gibt keine ungerade vollkommene Zahl.
\end{vermutung}

\noindent\textbf{Derzeit interessante offene Fragen umfassen:}
\begin{enumerate}[label=(\arabic*)]
  \item Kann $\omega(n) \geq 10$ bewiesen werden?
  \item Kann die Schranke $n > 10^{1500}$ wesentlich verbessert werden?
  \item Gibt es eine Primzahl $p$, die \emph{nicht} Euler-Primzahl einer
        ungeraden vollkommenen Zahl sein kann?
  \item Können quasivollkommene Zahlen ausgeschlossen werden?
  \item Gibt es ein algebraisch-zahlentheoretisches Argument, das in der
        Struktur $n = p^a m^2$ einen Widerspruch erzwingt?
  \item Könnten Siebbmethoden der analytischen Zahlentheorie einen
        nicht-rechnerischen Ansatz liefern?
\end{enumerate}

Das Problem berührt zutiefst die Struktur der multiplikativen Funktion $\sigma$
und Fragen über die Verteilung der Werte von $\sigma(n)/n$ nahe bei $2$ ---
einige der fundamentalsten ungelösten Fragen der elementaren Zahlentheorie.

% ============================================================
\begin{thebibliography}{99}

\bibitem{Euler1747}
L.\ Euler,
\textit{De numeris amicabilibus},
Nova Acta Eruditorum (1747);
auch in: Opera Omnia, Ser.\ I, Bd.\ 2, S.\ 86--162.

\bibitem{Touchard1953}
J.\ Touchard,
\textit{On prime numbers and perfect numbers},
Scripta Mathematica \textbf{19} (1953), 35--39.

\bibitem{Kanold1941}
H.-J.\ Kanold,
\textit{Untersuchungen über ungerade vollkommene Zahlen},
J.\ Reine Angew.\ Math.\ \textbf{183} (1941), 98--109.

\bibitem{Steuerwald1937}
R.\ Steuerwald,
\textit{Verschärfung eines Satzes über die Seite vollkommener Zahlen},
S.-B.\ math.-nat.\ Kl.\ Bayer.\ Akad.\ Wiss.\ 1937, 69--72.

\bibitem{Chein1979}
E.\ Z.\ Chein,
\textit{An odd perfect number has at least 8 prime factors},
Dissertation, Pennsylvania State University, 1979.

\bibitem{Hagis1980}
P.\ Hagis Jr.,
\textit{Outline of a proof that every odd perfect number has at least eight
prime factors},
Math.\ Comp.\ \textbf{35} (1980), 1027--1032.

\bibitem{BrentCohenRiele1991}
R.\ P.\ Brent, G.\ L.\ Cohen, H.\ J.\ J.\ te Riele,
\textit{Improved techniques for lower bounds for odd perfect numbers},
Math.\ Comp.\ \textbf{57} (1991), 857--868.

\bibitem{Iannucci1999}
D.\ E.\ Iannucci,
\textit{The third largest prime divisor of an odd perfect number exceeds
one hundred},
Math.\ Comp.\ \textbf{69} (2000), 867--879.

\bibitem{GotoOhno2008}
T.\ Goto und Y.\ Ohno,
\textit{Odd perfect numbers have a prime factor exceeding $10^8$},
Math.\ Comp.\ \textbf{77} (2008), 1859--1868.

\bibitem{LucaPomerance2010}
F.\ Luca und C.\ Pomerance,
\textit{On the radical of a perfect number},
New York J.\ Math.\ \textbf{16} (2010), 23--30.

\bibitem{OchemRao2012}
P.\ Ochem und M.\ Rao,
\textit{Odd perfect numbers are greater than $10^{1500}$},
Math.\ Comp.\ \textbf{81} (2012), 1869--1877.

\bibitem{OchemRao2014}
P.\ Ochem und M.\ Rao,
\textit{On the number of prime factors of an odd perfect number},
Math.\ Comp.\ \textbf{83} (2014), 2435--2439.

\bibitem{Nielsen2015}
P.\ P.\ Nielsen,
\textit{Odd perfect numbers, Diophantine equations, and upper bounds},
Math.\ Comp.\ \textbf{84} (2015), 2549--2567.

\bibitem{BroughanDelbourgoZhou}
K.\ A.\ Broughan, D.\ Delbourgo, Q.\ Zhou,
\textit{Improving the Chen and Chen result for odd perfect numbers},
Integers \textbf{13} (2013), A39.

\end{thebibliography}

\end{document}
